% Fontes locais no macOS e alternativas disponíveis no Overleaf.



\IfFontExistsTF{Arial}{
  \setmainfont{Arial}
  \setsansfont{Arial}
}{
  \setmainfont{TeX Gyre Heros}
  \setsansfont{TeX Gyre Heros}
}
\IfFontExistsTF{Menlo}{
  \setmonofont[Scale=0.86]{Menlo}
}{
  \setmonofont[Scale=0.86]{DejaVu Sans Mono}
}

\hypersetup{
  colorlinks=true,
  linkcolor=CorInstitucional,
  urlcolor=CorInstitucional,
  pdftitle={\TituloAvaliacao},
  pdfsubject={Modelo reutilizável de prova e gabarito},
  pdfauthor={\ElaboradorAvaliacao}
}

\setlength{\parindent}{0pt}
\setlength{\parskip}{0.35em}
\setlength{\arrayrulewidth}{0.75pt}
\setlist{nosep,leftmargin=1.6em}
\raggedbottom

% ----------------------------------------------------------------
% Rodapés
% ----------------------------------------------------------------
\pagestyle{fancy}
\fancyhf{}
\renewcommand{\headrulewidth}{0pt}
\renewcommand{\footrulewidth}{0pt}
\fancyfoot[L]{\footnotesize\sffamily\RodapeProva}
\fancyfoot[R]{\footnotesize\sffamily Página \thepage}

\newcommand{\usarRodapeProva}{%
  \fancyfoot[L]{\footnotesize\sffamily\RodapeProva}%
}

\newcommand{\usarRodapeGabarito}{%
  \fancyfoot[L]{\footnotesize\sffamily\RodapeGabarito}%
}

% ----------------------------------------------------------------
% Elementos do cabeçalho institucional
% ----------------------------------------------------------------
\newcommand{\logoavaliacao}{%
  \IfFileExists{\CaminhoLogoAvaliacaoPrincipal}{%
    \begin{minipage}{0.92\linewidth}
      \centering
      \includegraphics[width=0.96\linewidth,height=11mm,keepaspectratio]{\CaminhoLogoAvaliacaoPrincipal}\\[1.5mm]
      \IfFileExists{\CaminhoLogoAvaliacaoSecundario}{%
        \includegraphics[width=0.78\linewidth,height=9mm,keepaspectratio]{\CaminhoLogoAvaliacaoSecundario}%
      }{}%
    \end{minipage}%
  }{%
    \begingroup
    \setlength{\fboxsep}{4mm}%
    \colorbox{CorInstitucional!10}{\sffamily\bfseries\color{CorInstitucional}\MarcaAlternativaAvaliacao}%
    \endgroup
  }%
}

\newcommand{\campocabecalho}[2]{%
  \sffamily\textbf{#1}%
  \if\relax\detokenize{#2}\relax
  \else
    \newline\vspace{0.7mm}{\normalfont #2}%
  \fi
}

\usepackage{array}
\usepackage{tabularx}
\usepackage{multirow}
\usepackage{makecell} % Adicione este pacote no seu pré-ambulo para linhas mais espessas


\newcommand{\cabecalhoprova}{%
  \begingroup
  \setlength{\tabcolsep}{2mm}%
  \setlength{\arrayrulewidth}{1pt}%
  \renewcommand{\arraystretch}{1.25}%
  
  \setlength{\fboxrule}{2.5pt}%
  \setlength{\fboxsep}{0pt}%
  \fbox{%
    \begin{tcolorbox}[
      enhanced,
      sidebyside,
      sidebyside align=center,
      lefthand width=0.22\textwidth,
      sidebyside gap=0pt,
      valign=center,
      halign=center,
      boxrule=0pt,
      frame hidden,
      arc=0pt,
      colback=white,
      boxsep=0pt, % Zera o espaçamento interno do tcolorbox
      top=0pt, bottom=0pt, left=0pt, right=0pt,
      segmentation code={
        \draw[black, line width=1pt] 
          (segmentation.north) -- (segmentation.south);
      }
    ]
      % LADO ESQUERDO: Logo
      \logoavaliacao
      
      \tcblower
      
      % LADO DIREITO: Dados da avaliação
      \setlength{\tabcolsep}{2mm}%
      \setlength{\arrayrulewidth}{1pt}%
      \renewcommand{\arraystretch}{1.25}%
      \begin{tabularx}{\linewidth}{p{0.18\linewidth}|p{0.10\linewidth}|X|p{0.16\linewidth}|p{0.18\linewidth}}
        % Sem \hline no topo: a borda do \fbox já faz o fechamento superior
        \multicolumn{3}{l|}{\rule{0pt}{7mm}\campocabecalho{Nome:}{}} &
        \multicolumn{2}{l}{\campocabecalho{RA:}{}} \\
        
        \hline
        
        \campocabecalho{Data:}{\DataAvaliacao} &
        \campocabecalho{Turma:}{\TurmaAvaliacao} &
        \campocabecalho{Disciplina:}{\DisciplinaAvaliacao} &
        \campocabecalho{Docente:}{\DocenteAvaliacao} &
        \campocabecalho{Elaborador:}{\ElaboradorAvaliacao} \\
        
        \hline
        
        \multicolumn{5}{c}{\rule{0pt}{5.5mm}\sffamily\bfseries\Large\TituloAvaliacao \rule[-2.5mm]{0pt}{0pt}}
      \end{tabularx}%
    \end{tcolorbox}%
  }%
  \endgroup
}

\newcommand{\cabecalhogabarito}{%
  \begingroup
  \setlength{\tabcolsep}{2mm}%
  \setlength{\arrayrulewidth}{1pt}%
  \renewcommand{\arraystretch}{1.25}%
  
  \setlength{\fboxrule}{2.5pt}%
  \setlength{\fboxsep}{0pt}%
  \fbox{%
    \begin{tcolorbox}[
      enhanced,
      sidebyside,
      sidebyside align=center,
      lefthand width=0.22\textwidth,
      sidebyside gap=0pt,
      valign=center,
      halign=center,
      boxrule=0pt,
      frame hidden,
      arc=0pt,
      colback=white,
      boxsep=0pt, % Elimina folgas internas e linhas duplas nas bordas
      top=0pt, bottom=0pt, left=0pt, right=0pt,
      segmentation code={
        \draw[black, line width=1pt] 
          (segmentation.north) -- (segmentation.south);
      }
    ]
      % LADO ESQUERDO: Logo centralizada automaticamente no eixo vertical
      \logoavaliacao
      
      \tcblower
      
      % LADO DIREITO: Tabela do Gabarito
      \setlength{\tabcolsep}{2mm}%
      \setlength{\arrayrulewidth}{1pt}%
      \renewcommand{\arraystretch}{1.25}%
      \begin{tabularx}{\linewidth}{p{0.13\linewidth}|p{0.08\linewidth}|X|p{0.14\linewidth}|p{0.13\linewidth}}
        % Faixa superior destacada do GABARITO
        \multicolumn{5}{>{\columncolor{CorGabarito}}c}{\rule{0pt}{8mm}\sffamily\bfseries\fontsize{24}{28}\selectfont GABARITO \rule[-3mm]{0pt}{0pt}} \\
        
        \hline
        
        % Linha de informações intermediárias
        \campocabecalho{Semestre:}{\SemestreAvaliacao} &
        \campocabecalho{Turma:}{\TurmaAvaliacao} &
        \campocabecalho{Docente:}{\DocenteAvaliacao} &
        \campocabecalho{Disciplina:}{\DisciplinaAvaliacao} &
        \campocabecalho{Monitor:}{\MonitorAvaliacao} \\
        
        \hline
        
        % Título da avaliação
        \multicolumn{5}{c}{\rule{0pt}{5.5mm}\sffamily\bfseries\Large\TituloAvaliacao \rule[-2.5mm]{0pt}{0pt}}
      \end{tabularx}%
    \end{tcolorbox}%
  }%
  \endgroup
}
% ----------------------------------------------------------------
% Questões, alternativas e espaços de resposta
% ----------------------------------------------------------------
\newcounter{questao}
\newcounter{respostagabarito}

\newcommand{\pontuacao}[1]{%
  \if\relax\detokenize{#1}\relax\else\textbf{\textit{(#1)}}\fi
}

% Uso: \questao{1,0}{Enunciado} ou \questao{}{Enunciado sem nota}.
\newcommand{\questao}[2]{%
  \refstepcounter{questao}%
  \par\addvspace{0.85em}%
  \noindent\sffamily\textbf{\thequestao)}\quad\pontuacao{#1}\quad
  \normalfont\textit{#2}\par
}

% Uso: \subquestao{a}{0,75}{Texto da alternativa}.
\newcommand{\subquestao}[3]{%
  \par\addvspace{0.65em}%
  \noindent\hspace*{8mm}\sffamily\textbf{#1)}\quad\pontuacao{#2}\quad
  \normalfont #3\par
}

% Linhas abertas, úteis para respostas curtas.
\newcommand{\linhasresposta}[1]{%
  \par\begingroup
  \setlength{\parskip}{0pt}%
  \foreach \linha in {1,...,#1}{%
    \vspace{4mm}\noindent\rule{\linewidth}{0.38pt}\par
  }%
  \endgroup
}

% Caixa pautada e quebrável entre páginas, útil para respostas longas.
\usepackage{tcolorbox}
\tcbuselibrary{skins,breakable}

\newcommand{\respostalinhas}[1]{%
  \begin{tcolorbox}[
    enhanced,
    breakable,
    colback=white,
    frame hidden,
    boxrule=0pt,
    left=0mm, right=0mm, top=0mm, bottom=0mm,
    before skip=0.8em, after skip=0.8em,
    % 'overlay' aplica o desenho em todas as partes da caixa, quebrada ou não
    overlay={
      \begin{tcbclipinterior}
        \foreach \i in {1,...,#1} {
          \pgfmathsetmacro{\deslocamento}{\i * 22}
          \draw[black!28, line width=0.4pt] 
            ([yshift=-\deslocamento pt]frame.north west) -- 
            ([yshift=-\deslocamento pt]frame.north east);
        }
      \end{tcbclipinterior}
    }
  ]
  % Gera linhas com altura real para que o TeX saiba exatamente onde quebrar a página
  \begingroup
  \setlength{\parskip}{0pt}%
  \foreach \n in {1,...,#1}{%
    \strut\par\vspace{\dimexpr 22pt - \baselineskip\relax}%
  }%
  \endgroup
  \end{tcolorbox}%
}

\usepackage{tcolorbox}
\tcbuselibrary{skins, breakable}

\newcommand{\respostalinhacodigo}[1]{%
  \begin{tcolorbox}[
    enhanced,
    breakable,
    colback=white,
    frame hidden,
    boxrule=0pt,
    left=12mm, right=0mm, top=0mm, bottom=0mm,
    before skip=0.8em, after skip=0.8em,
    overlay={
      \begin{tcbclipinterior}
        % Calha de numeração (fundo cinza)
        \fill[black!4] (frame.north west) rectangle ([xshift=10mm]frame.south west);
        \draw[black!18, line width=0.5pt] 
          ([xshift=10mm]frame.north west) -- ([xshift=10mm]frame.south west);

        % Linhas pautadas horizontais
        \foreach \i in {1,...,#1} {
          \pgfmathsetmacro{\deslocamento}{\i * 22}
          \draw[black!20, line width=0.4pt] 
            ([yshift=-\deslocamento pt, xshift=10mm]frame.north west) -- 
            ([yshift=-\deslocamento pt]frame.north east);
        }
      \end{tcbclipinterior}
    }
  ]
  % Numeração sequencial com margem interna ajustada mais para a esquerda
  \begingroup
  \setlength{\parskip}{0pt}%
  \foreach \n in {1,...,#1}{%
    \noindent\hspace{-12mm}%
    \makebox[7.5mm][r]{\ttfamily\small\color{black!45}\n}\hspace{4.5mm}%
    \strut\par\vspace{\dimexpr 22pt - \baselineskip\relax}%
  }%
  \endgroup
  \end{tcolorbox}%
}

\usepackage{tcolorbox}
\tcbuselibrary{skins, breakable}

\newcommand{\caixalivre}[2]{%
  \begin{tcolorbox}[
    enhanced,
    breakable,
    colback=white,
    colframe=black!50, % Borda de tom neutro/sutil
    boxrule=0.6pt,
    arc=2pt, % Cantos ligeiramente arredondados
    left=3mm, right=3mm, top=2mm, bottom=2mm,
    before skip=0.8em, after skip=0.8em,
    attach title to upper,
    % Se um título for informado, formata no topo interno da caixa
    fonttitle=\sffamily\small\bfseries\color{black!60},
    title={#1\ifstrempty{#1}{}{\par\vspace{2mm}}}
  ]
  % Define a altura do bloco com base na quantidade de linhas desejada
  \vspace{\dimexpr 22pt * #2 - 2mm\relax}%
  \end{tcolorbox}%
}

\newcommand{\caixavf}[1]{%
  \begingroup
  \setlength{\tabcolsep}{3mm}%
  \renewcommand{\arraystretch}{1.3}%
  \begin{tabularx}{\textwidth}{c|X}
    \textbf{( V / F )} & \textbf{Afirmativa} \\ \hline
    #1
  \end{tabularx}
  \endgroup
}
% Uso:
% \caixavf{
%   [ \quad ] & O modelo em cascata é iterativo por definição. \\ \hline
%   [ \quad ] & O Scrum utiliza sprints de duração fixa.
% }

\newcommand{\caixarelacionar}[2]{%
  \begingroup
  \vspace{2mm}
  \noindent
  % Coluna A
  \begin{minipage}[t]{0.47\textwidth}
    \textbf{Coluna A}\\[1mm]
    \rule{\linewidth}{0.6pt}\\[2mm]
    \renewcommand{\arraystretch}{1.3}%
    \begin{tabular}{@{}l p{0.8\linewidth}@{}}
      #1
    \end{tabular}
  \end{minipage}%
  \hfill
  % Coluna B
  \begin{minipage}[t]{0.47\textwidth}
    \textbf{Coluna B}\\[1mm]
    \rule{\linewidth}{0.6pt}\\[2mm]
    \renewcommand{\arraystretch}{1.3}%
    \begin{tabular}{@{}l p{0.8\linewidth}@{}}
      #2
    \end{tabular}
  \end{minipage}
  \vspace{2mm}
  \endgroup
}

\usepackage{tcolorbox}
\tcbuselibrary{listings, skins}

\newtcblisting{codigolacuna}[1][]{%
  enhanced,
  colback=black!2,
  colframe=black!40,
  boxrule=0.5pt,
  arc=2pt,
  left=4mm, right=4mm, top=2mm, bottom=2mm,
  fontupper=\ttfamily\small,
  listing only,
  before skip=0.8em, after skip=0.8em,
  #1
}

% Macro auxiliar para desenhar a lacuna/traço no meio do código
\newcommand{\lacuna}[1][3cm]{\underline{\hspace{#1}}}

% Espaço em branco com altura livre: \espacoresposta{25mm}.
\newcommand{\espacoresposta}[1]{\par\vspace{#1}}

% Bloco de código neutro para questões técnicas.
\lstdefinestyle{codigoavaliacao}{
  basicstyle=\ttfamily\footnotesize,
  keywordstyle=\bfseries\color{CorInstitucional},
  commentstyle=\itshape\color{green!45!black},
  stringstyle=\color{red!65!black},
  numbers=left,
  numberstyle=\scriptsize\color{CorTextoAuxiliar},
  numbersep=8pt,
  showstringspaces=false,
  keepspaces=true,
  columns=fullflexible,
  breaklines=true,
  tabsize=2,
  upquote=true
}

\usepackage{tcolorbox}
\tcbuselibrary{listings, skins, breakable}

\newtcblisting{codigoquestao}[1][]{%
  enhanced,
  breakable,
  listing only,
  listing engine=listings,
  colback=CorCodigoAvaliacao, % Mantém sua cor de fundo original
  colframe=black!45,          % Borda neutra e elegante
  boxrule=0.5pt,
  arc=1.5mm,
  left=1mm, right=1.5mm, top=1.5mm, bottom=1.5mm,
  before skip=0.8em, after skip=0.85em,
  listing options={
    style=codigoavaliacao,
    numbers=left,                   % Ativa a numeração na esquerda
    numberstyle=\tiny\color{black!50}\ttfamily, % Estilo sutil do número
    numbersep=8pt,                  % Distância entre o número e o código
    firstnumber=1,
    stepnumber=1,
    #1                              % Permite passar parâmetros (ex: language=C)
  }
}

% Resposta numerada do gabarito. O segundo argumento pode ficar vazio.
\newcommand{\respostagabarito}[2][10mm]{%
  \refstepcounter{respostagabarito}%
  \par\addvspace{0.55em}%
  \noindent\sffamily\textbf{\therespostagabarito)}\quad\normalfont #2\par
  \vspace{#1}%
}
